\section{Linear Tropical Ideals}\label{sec:LinTropIdeals}
 

   This section will be focused on what we will call \textit{linear tropical ideals}, the tropical analogue of linear ideals. In the realizable case, linear ideals are perhaps the most simple and structured ideals of positive dimension, and any linear space can be realized as the variety of some linear ideal. However, when working with tropical ideals, valuated matroids, and their associated tropical linear spaces, it is already known that not all tropical linear spaces are tropically realizable varieties \cite{DR19}. Despite this complication in the theory, tropical linear spaces remain a reasonably well-behaved and well-studied class of tropical varieties, see \cite{Speyer08}, \cite{FR15}, \cite{GG18} for a few examples. In the remainder of this paper we will establish further combinatorial and geometric connections.
   
   A homogeneous ideal $I\subset k[x_1,\dots,x_n]$ is called linear if it is generated by linear forms or, equivalently, it is the ideal of polynomials vanishing on a linear space. While tropical ideals are generally not finitely generated, they do have a finite \textit{tropical basis} \cite{MR18}*{Theorem 5.9}, which is a set of polynomials in the ideal that cuts out the variety. This inspires the following definition.
   
    \begin{definition}    
    A homogeneous tropical ideal $I\subset \T[x_1,\dots, x_n]$ is a \textbf{linear tropical ideal} if its linear forms form a tropical basis, that is
    $$V(I)= V(I_1).$$
    \end{definition}
    The variety of a linear tropical ideal is necessarily a tropical linear space balanced with constant weights $1$.

    We say that an ideal $I\subset \T[x_1,\dots,x_n]$ is \textbf{saturated} with respect to a polynomial $f$ if $fp\in I$ implies $p\in I$. When $f=\prod_{i=1}^n x_i$ we will simply say that $I$ is saturated.
    
    \begin{remark}\label{rem:linSaturated}
    Given a linear tropical ideal $I\subset \T[x_1,\dots, x_n]$, we may express $I$ as a sequence of valuated matroids $(I_1,I_2,\dots,I_d,\dots)$. We will see that $I$ is saturated and that the truncation $(I_1,\dots,I_d)$ is a $d\th$ symmetric quasi-power of the matroid $I_1$.
    \end{remark}

\subsection{Connectivity}\label{subsec:connectivity}
 In order to prove the main results of this section, we will need to show that a $d$-dimensional tropical linear space has no proper tropical subvarieties of dimension $d$. To show this, we will prove a stronger result, which extends a valuable result in matroid theory to all valuated matroids. Specifically, we will show that the space of top-dimensional Minkowski weights of a tropical linear space is one-dimensional. We refer the reader to \cite{MS15}[Section 3.3] for relevant definitions and background for weighted polyhedral complexes, balancing, and the $\Star$ operation for a polyhedral complex.
 

If $\Sigma$ is a polyhedral complex in $\R^n$, we write $|\Sigma|$ to refer to the \textbf{support} of $\Sigma$, which is the set of points in $\Sigma$ without the polyhedral structure. Recall that the facet-ridge graph of $\Sigma$ is the graph whose vertices are given by the facets of $\Sigma$ and where there are edges between two vertices $\sigma_1$ and $\sigma_2$ if $\sigma_1$ and $\sigma_2$ intersect along a ridge of $\Sigma$. We say that $\Sigma$  is connected through codimension one if its facet-ridge graph $G(\Sigma)$ is a connected graph. We additionally say that $\Sigma$ is locally connected through codimension 1 if $G(\Star_w(\Sigma))$ is connected for every $w\in |\Sigma|$. 

\begin{proposition}\label{prop:connectedCodim1}
Let $\Sigma$ be a connected pure polyhedral complex that is locally connected through codimension one. Then $\Sigma$ is connected through codimension one.
\end{proposition}
\begin{proof}
    Let $\Sigma\subset \R^n$ be a connected, pure, polyhedral complex that is locally connected through codimension one. Further suppose towards a contradiction that $\Sigma$ is not connected through codimension one. Then the facet-ridge graph $G(\Sigma)$ has at least two connected components $C_1, C_2,\dots$. Since the vertices of $C_i$ are cones of $\Sigma$ we define the support of $C_i$ as
    $$|C_i|:= \bigcup_{\sigma\in V(C_i)} \sigma\subset |\Sigma|.$$
    Since $\Sigma$ is connected, there must exist $C_i$ and $C_j$ such that $|C_i|\cap |C_j|\neq \es$. Without loss of generality we suppose, suppose that $|C_1|\cap |C_2|\neq \es$, and let $w\in |C_1|\cap |C_2|$. Since $w\in |C_1|$ and $|C_2$ there is are facets  $\sigma_1\in V(C_1)$ and $\sigma_2\in V(C_2)$ such that $w\in \sigma_1$ and $\sigma_2$. Now, since $\Sigma$ is locally connected through codimension one, $G(\Star_w(\Sigma))$ is a connected graph. Moreover, by construction, $G(\Star_w(\Sigma))$ is a subgraph of $G(\Sigma)$. Since $\sigma_1$ and $\sigma_2$ are vertices of $\Star_w(\Sigma)$, there is a path connecting $\sigma_1$ and $\sigma_2$ in $G(\Star_w(\Sigma))$, however this now implies there is a path between $\sigma_1$ and $\sigma_2$ in $G(\Sigma)$. Therefore $C_1$ and $C_2$ are not distinct connected components, a contradiction. We conclude that $G(\Sigma)$ is necessarily a connected graph and therefore $\Sigma$ is connected through codimension one.
\end{proof}

The conditions of Proposition \ref{prop:connectedCodim1} are met by a large class of objects studied in tropical Geometry called tropical manifolds, which are meant to be a tropical analogue to classical manifolds. A $d$-dimensional \textbf{tropical manifold} is a polyhedral complex such that, for all $w\in |\Sigma|$, $\Star_w(\Sigma)$ is linearly isomorphic to the Bergman fan of some matroid that has dimension $d$. Of particular value to this paper is that every tropical linear space is a tropical manifold.

\begin{corollary}\label{cor:connectedCodim1}
    Any connected tropical manifold is connected through codimension one.
\end{corollary}
\begin{proof}
    The Bergman fan of any matroid is connected in codimension one. Since a connected tropical manifold $M$ is, by construction, locally the Bergman fan of a matroid, it is locally pure and locally connected in codimension one. This implies that $M$ is both pure and, by Proposition \ref{prop:connectedCodim1}, connected in codimension 1.
\end{proof}

\begin{definition}
    Let $\Sigma$ be a polyhedral complex. A $k\th$ Minkowski-weight of $\Sigma$ is an rational-valued function from the $k$-skeleton of $\Sigma$, denoted $\Sigma_k$, to the rationals
    \begin{align*}
        m:\Sigma_k \to \Q
    \end{align*}
    such that $\Sigma_k$ is a balanced polyhedral complex with weights given by $m$. 
\end{definition}


\begin{corollary}\label{cor:minkWeights}
The space of top-dimensional Minkowski-weights of a tropical manifold is a $1$-dimensional $\Q$-vector space and is generated by the constant weight function $w(\sigma)=1$. 
\end{corollary}
\begin{proof}
    Let $\Sigma$ be a $d$-dimensional tropical manifold. Let $w$ be a $d\th$ Minkowski weight of $\Sigma$ and suppose without loss of generality that $w(\sigma)=1$ for some facet $\sigma$. We claim that $w(\sigma')=1$ for all other facets of $\Sigma$. Notice that any facet $\sigma'$ adjacent to $\sigma$ in the facet-ridge graph $G(\Sigma)$ must have weight $1$. This follows as there exists a $w\in \sigma\cap\sigma'$ where $\star_w(\Sigma)=B(M_w)$ is the Bergman fan of some matroid $M_w$, which has $\sigma$ and $\sigma'$ as facets. Moreover, $B(M_w)$ is uniquely balanced with constant weights, and therefore since $w(\sigma)=1$ we must have $w(\sigma')=1$. By transitivity along the edges of $G(\Sigma)$, this implies that every facet in the same the connected component of $G(\Sigma)$ as $\sigma$ must have weight $1$. Since $G(\Sigma)$ is connected by Corollary \ref{cor:connectedCodim1} we see that every facet $\sigma'$ must have weight $w(\sigma')=1$.
    
\end{proof}
If $\Sigma'\subsetneq \Sigma$ is a containment of tropical varieties of equal dimension, then, after choosing a common refinement on their polyhedral structure, there is a natural embedding of the top-dimensional Minkowski weights of $\Sigma'$ into $\Sigma$. In particular, if the space of Minkowski weights of $\Sigma$ is one-dimensional, then $\Sigma$ can have no proper tropical subvarieties of equal dimension.

\begin{corollary}
    Connected tropical manifolds have no proper tropical subvarieties of equal dimension.
\end{corollary}

\subsection{Spaces of Symmetric Powers}\label{subsec:SpacesofSymPow}
Before continuing to prove that the theory of linear tropical ideals is equivalent to the theory of large symmetric powers of matroids, we must establish that the space of symmetric powers of a fixed matroid is a compact subspace. For this, we will need to define valuated matroids through their Pl\"ucker vectors.

\begin{definition}\label{def:valuatedMatroid}
   A \textbf{valuated matroid} $\nu$ is a vector in $\T^{\binom{[n]}{r}}$ that satisfies the following valuated exchange property. Let $B_1,B_2\in \binom{[n]}{r}$, then for every $e_1\in B_1\sm B_2$ there exists an $e_2\in B_2\sm B_1$ such that
    \begin{equation}\label{eqn:basisExchange}
    \nu(B_1)\odot\nu(B_2)\geq \nu(B_1\sm e_1\cup e_2)\odot\nu(B_2\sm e_2\cup e_2).
    \end{equation}
    The sets $B\subset [n]$ for which $\nu(B)<\infty$ are called \textbf{bases}. The bases of $\nu$ are the bases of a matroid that is called the \textbf{underlying matroid} of $\nu$, which we denote $\uMat(\nu).$
\end{definition}

The set of all Pl\"ucker vectors in $\TP^{\binom{n}{r}-1}$ is called the Dressian $\Dr(r,n)$, which is the moduli space of all valuated matroids of rank $r$ on $n$ elements. If the ground set $E$ of size $n$ is specified, we will write $\Dr(r,E)$ instead of $\Dr(r,n)$. Furthermore, we denote by $\Dr(M)$, the Dressian supported over $M$ or the Dressian of $M$, to be the set of valuated matroids whose underlying matroid is $M$. The Dressian $\Dr(r,n)$ is a closed polyhedral complex in $\TP^{\binom{n}{r}-1}$, moreover, since $\TP^{\binom{n}{r}-1}$ is compact, the Dressian is compact. However, $\Dr(M)$ is usually not compact, as it does not contain its boundary.

Let $M$ be a matroid with ground set $E$ and $T\subset E$. In order to specialize our focus to symmetric powers, we must understand how matroid restriction induces maps from $\Dr(M)$ to $\Dr(M|_T)$. Fix an independent set $I\subset E\sm T$ of rank $r-\rk(T)$ such that $T\cup I$ spans $E$. Given a valuated matroid $\nu\in \Dr(M)$, the restriction of $\nu$ to $T\subset E$ is given by
$$\nu|_T(B) = \nu(B\cup I).$$
The Pl\"ucker vector $\nu|_T\in \TP^{\binom{T}{\rk(T)}-1}$ does not depend on the choice of $I$ \cite{DW92}*{Proposition 1.2}. While the choice of $I$ does not matter when considering points of $\TP^{\binom{T}{\rk(T)}-1}$ we must consider the isometric embedding 
$$\imath_I:\TP^{\binom{T}{\rk(T)}-1}\hookrightarrow\TP^{\binom{[n]}{r}-1}$$ 
given by
 $$ (p_S)_{S\in\binom{T}{r}}\mapsto (a_B)_{B\in \binom{[n]}{r}} $$
 with
 $$a_B = \begin{cases}
 p_S & \text{if } B=S\cup I\\
 \infty & \text{if } B\neq S\cup I \text{ for any } S\subset [T].
 \end{cases}$$

 Given a Pl\"ucker vector $\nu\in \Dr(M)$, a restriction $\nu|_T\in \Dr(M|_{T})$, and a choice of $I$ with $T\cup I$ spanning $M$, we can interpret restriction as a coordinate projection 
 $$\pi_{T\cup I}: \TP^{\binom{n}{r}}\to \TP^{\binom{n}{r}}$$
 
 given by
$$\nu\mapsto \imath_I(\nu|_T).$$


\begin{proposition}\label{prop:plConvergence}
Let $\set{\nu_i}_{i=1}^\infty$ be a convergent sequence of Pl\"ucker vectors in $\Dr(r,n)\subset \TP^{\binom{n}{r}-1}$ with limit $\nu$. Then for any $T\subset [n]$ where $\nu|_T$ has positive rank, the sequence $\set{\nu_i|_T}_{i=1}^\infty$ converges to $\nu|_T$ in $\Dr(\rk(T), T)$
\end{proposition}
\begin{proof}
 Since $\Dr(r,n)$ is a finite union of the sets $\Dr(M)$, we may assume that the entire sequence $\set{\nu_i}_{i=1}^\infty$ lies in $\Dr(M)$ for fixed $M$, but possibly converges to a point in the boundary of $\Dr(M)$.  Since $\nu$ has at least one basis $B$, we choose to present our vectors such that $(\nu_i)_B = 0$. 
 Consider that for each $A\in \binom{[n]}{r}$ that 
 $$\lim_{i\to \infty} (\nu_i)_A = \nu_A.$$
 For any choice of independent set $I\subset [n]\sm T$ of rank $r-\rk(T)$ such that $T\cup I$ spans $M$ we have the following. Let $S\in \binom{T}{\rk(T)}$, the sequence
 $$\set{(\nu_i|_{T})_S}_{i=1}^\infty$$ 
 converges to $\nu_{S\cup I}$ as
 $$\lim_{i\to\infty}\set{(\nu_i|_{T})_S}_{i=1}^\infty=\lim_{i\to\infty}\set{(\nu_i)_{S\cup I}}_{i=1}^\infty=\nu_{S\cup I}.$$
 
Let $\rho\in \TP^{\binom{T}{\rk(T)}}$ be the vector with coordinates 
$$\rho_S = \nu_{S\cup I}.$$
We prove that there is a choice of independent set $I$ for which $T\cup I$ spans $\uMat(\nu)$ and that $I$ has rank $r-\rk(T)$ complementary rank in $\uMat(\nu)$, as given this we then have that $\rho_S = \imath\inv(\pi_{T\cup I}(\nu))= \nu|_T.$ Since $\rho_S = \nu_{S\cup I}$, we see that $\rho_S<\infty$ if and only if $S\cup I$ is a basis of $\uMat(\nu)$. Since $\rho$ has at least one basis then $\rho_S = \nu_{S\cup I}$ and $S\cup I$ is a basis of $\uMat(\nu)$. In particular $S\cup I\subset T\cup I$ spans $\uMat(\nu)$, $S$ has rank $\rk(T)$, and by construction $I$ rank $r-\rk(T)$, as desired. Therefore $\rho = \imath\inv(\pi_{T\cup I}(\nu)) = \nu|_T$, as desired.
\end{proof}



\begin{proposition}\label{prop:closedcompact}
Let $M$ be a fixed valuated matroid of rank $r$ in $\T^{[n]}$. Set
$$r_d:= \binom{r+d-1}{d}$$
$$n_d:= \binom{n+d-1}{d}.$$
Then the collection $\Dr_d(M)$ of $d\th$ symmetric powers of $M$ form a compact polyhedral complex in $\Dr(r_d,n_d)$.
\end{proposition}
\begin{proof}
Fix a valuated matroid $M$. We note that $M_d\in \Dr_d(M)$ determines $(M,M_2,\dots,M_d)$ uniquely, thus we only need to record the Pl\"ucker vector of $M_d$ in this construction, hence we work in $\TP^{\binom{n_d}{r_d}-1}$. 

It is sufficient to prove that the set of vectors in $\TP^{\binom{n_d}{r_d}-1}$ satisfying Isomorphism \ref{eqn:symprodEmbedding} forms a closed polyhedral complex.

Let $p$ be the Pl\"ucker vector of $M_{d-1}$ and for each $\omega\in [n]$ fix a set $I_\omega$ such that $\omega\Sym_{d-1}(E)\cup I_\omega$ spans $\uMat(M_d)$. For simplicity, we abbreviate the associated coordinate projection by
$$\pi_\omega: = \pi|_{(\omega \Sym_{d-1}(E)\cup I_\omega)}.$$

The set of vectors in $\Dr(r_d,n_d)$ that satisfy Isomorphism \ref{eqn:symprodEmbedding} for the embedding of $M_{d-1}$ by $\omega$ into $M_d$, for all $\omega\in [n]$, are those vectors in the intersection of the $Dr(r_d,n_d)$ and the following set
$$S:=\bigcap_{\omega\in [n]}\set{v\in \TP^{\binom{n_d}{r_d}-1}: \pi_\omega(v) = p}.$$

We now note that the set
$$S_\omega:= \set{v\in \TP^{\binom{n_d}{r_d}-1}: \pi_{\omega}(v) = p}$$

is a closed polyhedral subset of $\TP^{\binom{n_d}{r_d}}$, and therefore the finite intersection $S$ over all $S_\omega$ is closed and polyhedral. We conclude that $\Dr_d(M)=S\cap \Dr(r_d,n_d)$ is a compact polyhedral complex. 
\end{proof}

\subsection{Hilbert Functions}\label{subsec:HF}
    Now that we've established that a tropical linear space cannot have any proper tropical subvarieties of equal dimension, we can return to our discussion of linear tropical ideals. Recall that the \textbf{Hilbert function} of a homogeneous tropical ideal $I\subset \T[x_1,\dots,x_n]$ is given by
    $$H_I(d):= \rk(I_d).$$
    While this definition is given for homogeneous tropical ideals, it extends to a definition for ideals in $\T[x_0,\dots x_n]$ or $\T[x_1^{\pm1},\dots x_n^{\pm1}]$. See \cite{MR20}*{Lemma 2.1} for a comprehensive description. 
    
    The Hilbert function of a tropical ideal $I$ tells us valuable information about the dimension of the tropical variety $V(I)$, and this fact motivates what follows. In order to establish bounds on the value $H_I(d)$ we will need to define the independence complex of a tropcial ideal $I$.
\begin{definition}\label{def:independenceComplex}
    The \textbf{independence complex} of a tropical ideal $I\subset\T[x_1,\dots,x_n]$ is the set system on $[n]$ given by
    $$\sI(I):= \set{A\subset [n]: I\cap \T[x_i]_{i\in A} = \set{\infty}}$$
    or equivalently
    $$\sI(I):= \set{A\subset [n]: \pi_{A}(V(I))=\T^{A} }.$$
\end{definition}

    
\begin{proposition}\label{prop:linearHFequiv}
    Let $I\in \T[x_1,\dots,x_n]$ be a tropical ideal with $\rk(I_1)=r$. Then $I$ is linear if and only if
    $$H_I(d)=\binom{r+d-1}{d}.$$
\end{proposition}
To avoid confusion in the following proof, we note that our convention of $H_I(d)$ and that used in the paper \cite{DR19} differ. In this paper $H_I(d)=\rk(I_d)$, whereas in \cite{DR19} the convention $H_I(d)=\rk(I_{\leq d})$ is used.
\begin{proof}
Let $I$ be a linear tropical ideal with $\rk(I_1)=r$ and let $M$ be the underlying matroid $\uMat(I_1)$. Then by definition $V(I)=V(I_1)$ and $H_I(1)=\rk(I_1)=r$. By \cite{DR19}*{Proposition 5.5} we can immediately conclude that
$$H_I(d)\leq \binom{r+d-1}{d}.$$

To impose a lower bound, we notice that $\pi_A(V(I))=\T^A$ if and only if $\pi_A(B(M))=\T^A$. In particular, the independence complex $\sI(I_1)$ is equal to that of $B(M)$. \cite{Yu17}*{Lemma 3} then provides that $\sI(I_1)$ is the collection of independent sets of $M$, and therefore $\sI(I_1)$ contains a set of size $r$. By \cite{DR19}*{Proposition 5.4} we may conclude that
$$H_I(d)\geq \binom{r+d-1}{d}.$$ 
Combining these two inequalities provides that $H_I(d)= \binom{r+d-1}{d}$, as desired.

Conversely, suppose that $I$ has Hilbert function exactly $H_I(d)=\binom{r+d-1}{d}$. Then $I_1$ is a valuated matroid of rank $H_I(1)=r$ and $V(I_1)$ is a tropical linear space of dimension $r-1$. Moreover, the dimension of $V(I)$ is equal to the degree of the Hilbert polynomial of $H_I(d)$ \cite{MR20}*{Theorem 4.3}, which is a degree $r-1$ polynomial. We now note that $V(I)$ is a tropical subvariety of $V(I_1)$ which, by Corollary \ref{cor:minkWeights} contains no proper tropical subvariety of dimension $r-1$. Thus we must have $V(I)=V(I_1)$, and therefore $I$ is linear.
\end{proof}

\begin{example}\label{ex:nonrealizableSymprod}
Consider the non-realizable tropical ideal $I$ from \cite{MR16}*{Example 2.8}. In \cite{MR16}*{Example 5.15} it is shown that $I$ has variety equal to the tropical line. In particular, $I_1$ is a tropical basis of $I$ and so $I$ is tropically linear. Since $I$ is a saturated ideal, we see that $I_3$ is a third symmetric power of $I_1=U_{2,3}$ that is not realizable as a linear matroid over any field.
\end{example}

\begin{corollary}\label{cor:LTIsaturated}
If $I\subset \T[x_1,\dots,x_n]$ is a linear tropical ideal then

$$I = (\T[x_1^{\pm1},\dots,x_n^{\pm1}]\cdot I )\cap \T[x_1,\dots,x_n].$$
In other words, $I$ is saturated with respect to $\prod_{i=1}^n x_i$. 
\end{corollary}
\begin{proof}
Let $I$ be a linear tropical ideal, but suppose that towards a contradiction that $I$ is not saturated. Let $I_d$ be the minimal degree that is not saturated and let $f= x^\omega \cdot g\in I_d$ where $x^\omega$ is a degree $i$ monomial, $g$ is not divisible by any monomial, and $g\notin I_{d-i}$. This implies $x^\omega\cdot I_{d-i}\subsetneq I_d|_{x^\omega}$ are nested tropical linear spaces. In particular $I_d|_{x^\omega}$ is not full rank and therefore $$V(I_d|_{x^\omega})\subsetneq B(M).$$
However this immediately implies $V(I)\subsetneq V(I_1)$, a contradiction. Thus $I$ must be saturated.
\end{proof}

\begin{remark}\label{rem:bfCharacterisation}
The proofs of \cite{DR19}[Proposition 5.4, Proposition 5.5], which are used to prove Proposition \ref{prop:linearHFequiv}, are made at the level of tropical linear spaces, and not tropical ideals. In particular, they can be restricted to any $d\th$ symmetric power of a matroid. Indeed, we may conclude $(M_1,\dots,M_d)$ is a symmetric power of $M$ if and only if $V(M_d)= B(M)$. 
\end{remark}

\begin{corollary}
Statements (1) and (4) are equivalent in Proposition \ref{prop:symprodCharacterisations}
\end{corollary}

\begin{corollary}\label{cor:hfSymprodEquiv}
Let $I\subset \T[x_1,\dots,x_n]$ be a tropical ideal. Then $I$ is linear if and only if, for every $d\in \N$, $(I_1,\dots,I_d)$ is a $d\th$ symmetric power of $I_1$. 
\end{corollary}
\begin{proof}
    Corollary \ref{cor:LTIsaturated} and Proposition \ref{prop:linearHFequiv} imply that if $I\subset \T[x_1,\dots,x_n]$ is a linear tropical ideal then the sequences $(I_1,\dots,I_d)$ are $d\th$ symmetric powers of $I_1$. Moreover, if $I$ is a tropical ideal such that $(I_1,\dots,I_d)$ are $d\th$ symmetric powers of $I_1$, then 
    $$V(I) = \cap_{d=1}^\infty V(I_d),$$
    where, by Proposition \ref{prop:symprodCharacterisations}, we have $V(I_d)=V(I_1)$. We conclude that
    $$V(I)=\bigcap_{d=1}^\infty V(I_1) = V(I_1),$$
    so $I$ is linear.
\end{proof}

 We will use these result to understand how matroid symmetric powers tie in to the theory tropical ideals in a very natural way.

\subsection{Tropically Realizable Matroids}\label{subsec:TropRealizableMatroids}

Given a linear tropical ideal $I$, the data of its variety is encoded entirely in $I_1$, which is itself a valuated matroid, and we have seen that $I$ necessarily encodes infinitely large symmetric powers of $I_1$. We will now prove that these are in fact equivalent theories.

Recall that a tropical variety $X$ is said to be tropically realizable if there is a tropical ideal $I$ for which $V(I)= X$. We extended this terminology to matroids by saying a valuated matroid $\nu$ on $[n]$ is tropically realizable if there exists a linear tropical ideal $I$ in $\T[x_1,\dots,x_n]$ for which $V(I)=\trop(\nu)$, where $\trop(\nu)$ is the tropical linear space of $\nu$.

\begin{theorem}\label{thm:linearMatroidCorres}
Let $\nu$ be a valuated matroid. Then $\nu$ is tropically realizable if and only if for every $d$ some symmetric power $(\nu_1,\dots,\nu_d)$ exists.
\end{theorem}

\begin{proof}
Since Proposition \ref{prop:plConvergence} requires $\nu$ to have positive rank we note that if $\nu$ is rank $0$ the proof is trivial as rank $0$ valuated matroids have unique symmetric powers. Thus we may assume $\nu$ has positive rank in the remainder of the proof.

Let $\nu$ be a tropically realizable matroid. Then there exists a tropical ideal $I\subset \T[x_1,\dots, x_n]$ with $V(I)=\trop(\nu)$ and $H_I(d)=\binom{\rk(I_1)+d-1}{d}$. Corollary \ref{cor:hfSymprodEquiv} now provides that $(I_1,\dots,I_d)$ is a $d\th$ symmetric power of $I_1=\nu$ for all $d\in \N$. This proves the implication.

For the converse, suppose that we have an infinite collection of symmetric powers of $\nu$ that become arbitrarily large, say 
$$S = \set{(\nu=\nu_{1,d},\nu_{2,d},\dots, \nu_{d,d})}_{d=1}^\infty.$$
We claim that the existence of $S$ implies that there is an infinite sequence $J$ of symmetric powers of $\nu$
$$J = (\nu, \nu_2, \dots, \nu_d, \dots)$$
where now the matroids $\nu_i$ are fixed between all $d\th$ symmetric powers given in $J$. Given that $J$ exists we will see that $J$ defines a tropical ideal via the circuit polynomials of each $\nu_d$. Thus we construct $J$ given $S$.

Let $S_j$ be any subset of $S$. We denote by $S_{j,d}$ the multiset of Pl\"ucker vectors
$$\set{\rho_d\in \Dr_d(\nu): (\rho_1,\rho_2,\dots, \rho_d,\dots, \rho_k)\in S_j}.$$

Let $\nu_1 = \nu$ and $S_1=S$. We recursively define the valuated matroid $\nu_d$ and the subset $S_d\subset S$ as follows. Let $\nu_d$ be an accumulation point or any point of infinite multiplicity in $S_{d-1,d}$. Let $S_d$ be a subset of $S_{d-1}$ for which $\nu_d$ is the unique accumulation point or point of infinite multiplicity in $S_{d,d}$. We claim that for every $d\in \N$ that $\nu_d$ and $S_d$ exist and that $(\nu_1,\dots \nu_d)$ is a $d\th$ symmetric power of $\nu$.

Proposition \ref{prop:closedcompact} provides that since $S_{d-1,d}$ is an infinite multiset contained in $\Dr_{d-1}(\nu)$ that $S_{d-1,d}$ must have an accumulation point or a point of infinite multiplicity. Fix $\nu_d$ to be any such point. By construction of $\nu_d$, there is necessarily an infinite subset $S_d\subset S_{d-1}$ such that $\nu_d$ is the unique accumulation point in $S_{d,d}$. 

We now claim that $(\nu_1,\dots,\nu_d)$ is a $d\th$ symmetric power of $\nu$. Let 
$$T:=\set{\rho_{d,i}}_{i=1}^\infty$$
be a convergent sequence of elements of $S_{d,d}$. By construction $T$ necessarily converges to $\nu_d$. Proposition \ref{prop:plConvergence} then implies that for any $e \in E(\nu)$ that the sequence
$$T|_e:=\set{\rho_{d,i}|_{e}}_{i=1}^\infty$$
converges to $\nu_d|_e$. Since each $\rho_{d,i}$ is a $d\th$ symmetric power of $\nu$, Isomorphism \ref{eqn:symprodEmbedding} provides that we have an equality of Pl\"ucker vectors. $\rho_{d,i}|_{e}=\rho_{d-1,i}$.
In particular, the sequence
$$T|_e = \set{\rho_{d-1,i}}_{i=1}^\infty\subset S_{d-1,d-1}.$$
By construction, any convergent sequence in $S_{d-1,d-1}$ necessarily converges to $\nu_{d-1}$. Thus $\nu_d|_e=\nu_{d-1}$. This implies that $(\nu_1,\dots,\nu_{d-1},\nu_d)$ is a symmetric power of $\nu_1=\nu$ if and only if $(\nu_1,\dots,\nu_{d-1})$ is a symmetric power. In the case where $d=2$ this follows immediately, and therefore holds for all $d$, as desired. We conclude that the infinite sequence $J=(\nu_1,\dots,\nu_d, \dots)$ exists.

We claim that for any sequence $J$ obtained by the above construction for any choice of accumulation points or points of infinite multiplicity in each step, that $J$ provides a linear tropical ideal. Consider that Isomorphism \ref{eqn:symprodEmbedding} provides that the polynomials associated to each $\nu_d$ defined by $J$ form a saturated tropical ideal $I\subset \T[x_1,\dots,x_n]$ with Hilbert function $\binom{\rk(\nu)+d-1}{d}$. Proposition \ref{prop:linearHFequiv} now provides that $I$ is a linear tropical ideal with $V(I)=V(I_1)=\trop(\nu)$, as desired.
\end{proof}

\begin{remark}
It remains an open problem whether or not there is a tropical linear space that is tropically realizable but not realizable over a field. 
\end{remark}

\subsection{Minor Closed}\label{subsec:MinorClosed}
    In this section we will show how tropical geometry allows us to easily prove that the class of tropically realizable matroids is minor closed and that the class of matroids with $d\th$ symmetric power is minor closed.
    
\begin{theorem}\label{thm:minorClosed}
The class of tropically realizable matroids is minor closed. Specifically, given a tropical realization $I\subset \T[x_1,\dots,x_n]$ of the valuated matroid $\nu$, the contraction $\nu/n$ is tropically realized by the specialization $I_{x_n=\infty}$, and the deletion $\nu\sm n$ is tropically realized by $I\cap\T[x_1,\dots,x_{n-1}]$.
\end{theorem}
\begin{proof}
Without loss of generality, consider the ideal $J:= I|_{x_n=\infty}$. We claim $J$ tropically realizes $\nu/n$. We first prove that $V(J) = V(J_1)=\nu/n$. Note that
$$V(J)= \pi_{[n-1]}(V(I)\cap \set{x: x_n = \infty}) = \pi_{[n-1]}(\trop(\nu)\cap\set{x: x_n=\infty}).$$
To see that $V(J)\subset \trop(\nu/n)$ it is sufficient to notice that $J_1 = \nu/n$. This holds as $I_1=\nu$ and $J_1=\nu/n$ by construction. Therefore $V(J)\subset V(J_1)\subset \trop(\nu/n)$. 

Now let $\omega\in \trop(\nu/n)$. We must show that every polynomial in $J$ attains its minimum twice on $\omega$. Note that $\omega\in \trop(\nu/n)$ if and only if $\omega'=(\omega,\infty)$ is in $\trop(\nu)=V(I)$. We now show that this implies $\omega\in V(J)$. Consider that every polynomial $g$ of $J$ is of the form $f|_{x_n=\infty}$ for some polynomial $f\in I$. In particular $f\in I$ attains its minimum twice or equals $\infty$ at $\omega'$ if and only if $g$ attains its minimum twice or equals $\infty$ at $\omega$. Since $\omega'\in V(I)$ it follows that $\omega\in V(J)$, as desired. 

That $J_d$ is a symmetric quasi power of $J_1$ is a consequence of Corollary \ref{cor:LTIsaturated}.

\commentout{as otherwise, for some degree $i$ monomial $x^{\bfu}$ and some degree $d$, $J_d|_{\bfu}$ would be a nontrivial quotient of $J_{d-i}$. In particular, the rank of $J_d|_{\bfu}$ would be strictly less than that of $J_{d-i}$ and Proposition \ref{prop:symprodCharacterisations} would imply $V(J)\subset V(J_d|_{\bfu}\subset\neq \trop(\nu/n)$, a contradiction.} 
As for the matroid $M\sm n$, consider the tropical ideal $J := I\cap \T[x_1,\dots,x_{n-1}]$. The variety of $J$ is $\pi_{[n-1]}(B(M))$ by \cite{MR20}*{Theorem 4.7}, which is equal to $B(M\sm i)$, \cite{dMRS20}*{Section 5.1}. By construction $I_1$ is equal to $M\sm i$. Therefore $I$ tropically realizes $M$. 
\end{proof}


\begin{lemma}\label{lem:geometricDeletion}
    If $\nu$ is a valuated matroid on $[n]$ and $\nu_d$ is a $d\th$ symmetric power of $\nu$. Then the valuated matroid
    $$P:=\nu_d\sm(n\cdot \Sym_{d-1}([n]))$$
    satisfies
    $$V(P) = \pi_{[n-1]}(\trop(\nu)).$$
    In particular, $P$ is a $d\th$ symmetric power of $\nu\sm n$.
\end{lemma}
\begin{proof}
Proposition \ref{prop:symprodCharacterisations} provides that $V(\nu_d)=\trop(\nu)$.

Let $w\in V(P)$. Then every circuit polynomial $f$ of $P$ attains its minimum twice or is equal to $\infty$ at $w$. However, this implies that every circuit polynomial of $\nu_d$ attains its minimum twice or is equal to infinity at the point $w'=(w,\infty).$ In particular, $w'\in V(\nu_d)$ and $w=\pi_{[n-1]}(w')\in \pi_{[n-1]}(V(\nu_d))=\pi_{[n-1]}(\trop(\nu)).$

Now let $w\in \pi_{[n-1]}(V(\nu_d))$. Then there is a $w' \in V(\nu_d)$ with $\pi_{[n-1]}(w')=w$. In particular, every circuit polynomial of $\nu_d$ attains its minimum twice or is equal to $\infty$ at $w'$. However, we note a circuit polynomial $f$ of $P$ is also circuit polynomial of $\nu_d$ and that $f$ attains its minimum twice or is equal to $\infty$ at $w'$ if and only if $f$ attains its minimum twice or is equal to $\infty$ at $w$. We conclude that $w\in V(P)$.

Since $\pi_{[n-1]}(\trop(\nu))=\trop(\nu\sm n)$ we conclude that $V(P)=\trop(\nu\sm n)$. By construction, $P$ is a symmetric quasi power of $\nu\sm n$ and Proposition \ref{prop:symprodCharacterisations} then implies that $P$ is a symmetric power of $\nu\sm n$, as desired.
\end{proof}

\begin{corollary}\label{cor:minorClosed}
    The class of matroids with $d\th$ symmetric power is minor closed.
\end{corollary}
\begin{proof}
    The result follows from the argument used in Theorem \ref{thm:minorClosed}, replacing \cite{MR20}*{Theorem 4.7} with Lemma \ref{lem:geometricDeletion}.
\end{proof}

\commentout{
While it is valuable to know the class of matroids with $d\th$ symmetric power is minor closed, the class of matroids with any symmetric powers at all remains elusive. However, we can deduce a structural result that provides insight into how matroid symmetric powers are built. 

Let $\nu$ be a valuated matroid with symmetric power $(\nu_1,\dots,\nu_d)$. By Theorem \ref{thm:minorClosed}, for any subset $S$ of the ground set, the matroid $\nu_d/(S\cdot \Sym_{d-1}(E))$ is a symmetric power of $\nu/S$, thus we define

$$(\nu/S)_d:=\nu_d/(S\cdot \Sym_{d-1}(E)).$$


\begin{corollary}
    Let $\nu$ be a valuated matroid on the ground set $E$ and let $(\nu_1,\dots,\nu_d)$ be a $d\th$ symmetric power of $\nu$. Then for $S\subset E$ and $a\notin S$ we have
    $$(\nu/S)_d|_a = \nu_d|_a/a(S\cdot \Sym_{d-2}(E\sm a))\cong (\nu/S_d).$$
    
\end{corollary}
\begin{proof}
Let 
$$P_1=\nu_d|_a/a(S\cdot \Sym_{d-2}(E\sm a))$$
and
$$P_2=\nu_d/(S\cdot \Sym_{d}(E))|_a.$$
Considering the vectors of $P_1$ and $P_2$ we notice that $P_2$ is a quotient of $P_1$ as every vector of $P_1$ must be a vector of $P_2$. By Theorem \ref{thm:minorClosed} we notice that both $P_1$ and $P_2$ are $d-1$ symmetric powers of $M/S$. In particular $\rk(P_1)=\rk(P_2)$. Since $P_2$ is a quotient of $P_1$ of equal rank, this implies that $P_1=P_2$, as desired.
\end{proof}}
%%%%%%%%%%%%%%%%%%%%%%%%%%%%%%%%%%%%%%%%%%%%%%%%%%%%

%%%%%%%%%%%%%%%%%%%%%%%%%%%%%%%%%%%%%%%%%%%%%%%%%%%%

%%%%%%%%%%%%%%%%%%%%%%%%%%%%%%%%%%%%%%%%%%%%%%%%%%%%
\subsection{Tropically Algebraic and Set-realizable Matroids}\label{subsec:A&StRealizable}

    Matroids appear in many ways in the study of tropical ideals. In the previous sections we considered only linear tropical ideals, whose variety is a tropical linear space with weights $1$. However, we can just as easily ask about tropical linear spaces balanced with weights other than $1$. This remains largely an unexplored avenue of research. A valuated matroid $\nu$ is called \textbf{tropically set-realizable} if $\trop(\nu)$ is the support of $V(I)$ for some tropical ideal $I$. That is, we consider equality of $\trop(\nu)$ with $V(I)$ as sets with no polyhedral structure. See \cite{Yu17} for initial work in this direction. 
    
    Stemming from tropical set-realizability is also the study of \textbf{tropically algebraic matroids} --- matroids that appear as the independence complex of a tropical ideal, as in Definition \ref{def:independenceComplex}. Tropically algebraic matroids were first discussed in \cite{DR19} but there is yet any work exploring this topic.
    
    Restricting to the case of matroids rather than valuated matroids, we observe that the classes of tropically realizable, tropically set-realizable, and tropically algebraic matroids form nested classes of matroids in increasing order. There are yet any proofs that any of these containment relations is strict. However, in this section we conclude that all three classes are minor closed.

\begin{proposition}
The class of tropically set-realizable valuated matroids is minor closed.
\end{proposition}
\begin{proof}
The exact proof of Theorem \ref{thm:minorClosed} applies here without the requirement that $I_1=M$, as the weights associated to the maximal cones are no longer considered.
\end{proof}

While minor-closedness of tropically set-realisable matroids follows the same proof as used for tropically realizable matroids, tropically algebraic matroids require a slightly different approach.

If the independence complex of $I$ forms the independent sets of a matroid $M$ we say that $I$ is a matroidal tropical ideal. Furthermore, if there is a tropical ideal with $\sI(I)=M$ then we say that $M$ is tropically algebraic. Algebraic matroids are tropically algebraic \cite{Yu17}*{Lemma 3}, but it is still open whether or not there is a tropically algebraic matroid that is not algebraic.

\begin{proposition}\label{prop:algMinorClosed}
The class of tropically algebraic matroids is minor closed.
\end{proposition}
\begin{proof}
Suppose $M$ is a tropically algebraic matroid, and let $I$ be a tropical ideal whose independence complex is $M$. Let without loss of generality we delete $n$. We claim that $J:=I\cap \T[x_1,\dots,x_{n-1}]$ has $\mathcal I(J) = M\sm n$. To see this, we first note \cite{MR20}*{Theorem 4.7} provides that
$$V(J)= \pi_{[n-1]}(V(I)).$$

A set $T$ is independent in $M\sm n$ if and only if $T$ is independent in $M$ and does not contain $n$. Thus we must show that $\pi_T(V(J))=\T^T$ if and only if $\pi_T(V(I))=\T^T$. To see this, we need only observe that $\pi_T(V(J))=\pi_T(\pi_{[n-1]}(V(I)))=\pi_T(V(I))$, as desired.

To prove contraction, if $n$ is a loop of $M$, then $M/n =M\sm n$, so we choose to delete $n$ instead and do the above. Otherwise $n$ is not a loop and we do the following. Consider the specialization $J:= I|_{x_n=0}$. The independence complex of $J$ is 
$$\mathcal I(J) = \set{T\subset [n-1] | J \cap \R[\set{x_i}_{i\in T}] = \set{\infty}}.$$

We claim $\mathcal I(J) = \mathcal I(M/n)$. In particular, we must show that $T\in \mathcal I(J)$ if and only if $T\cup i$ is independent in $\mathcal I(I)$. Since $\sI(J)$ is a simplicial complex, it is sufficient to prove this for its bases.


First let $B$ be a maximal subset of $\mathcal I(J)$. We claim $B\cup n$ is in $\mathcal I(I)$. Suppose not. Then there is a nonzero polynomial $f$ in $I\cap \T[x_i: i \in B\cup n]$. This implies that $f|_{x_n=0}\neq \infty$, where $f|_{x_n=0}\in J$, a contradiction. Thus $\mathcal I(J)\subset \mathcal I(M/n)$

That $\mathcal I(M/n)\subset \mathcal I(J)$ is immediate, as if $B\in \mathcal I(M/n)$ is a basis, we have that $B\cup \in \mathcal I(M)=\mathcal I(I)$, in particular, $I\cap \T[x_i: i\in (B\cup n)]=\set{\infty}$, and thus $J\cap \T[x_i: i\in B]=\set{\infty}$. This concludes the proof.
\end{proof}

In the proof of Proposition \ref{prop:algMinorClosed}, we specialize our ideals at $x_n=0$, rather than $x_n=\infty$. However, in the case where $V(I)=\trop(\nu)$ for some valuated matroid $\nu$, both specialization at $x_n=0$ and $x_n=\infty$ yield tropical ideals that tropically algebraically represent $M/n$.


%%%%%%%%%%%%%%%%%%%%%%%%%%%%%%%%%%%%%%%%%%%%%%%%

%%%%%%%%%%%%%%%%%%%%%%%%%%%%%%%%%%%%%%%%%%%%%%%%

%%%%%%%%%%%%%%%%%%%%%%%%%%%%%%%%%%%%%%%%%%%%%%%%

%%%%%%%%%%%%%%%%%%%%%%%%%%%%%%%%%%%%%%%%%%%%%%%%
\commentout{
\begin{proposition}\label{prop:plConvergence}
Let $\set{\nu_i}_{i=1}^\infty$ be a convergent sequence of Pl\"ucker vectors in $\Dr(M)$ with limit $\nu$. Then for any $T\subset [n]$, the sequence $\set{\nu_i|_T}_{i=1}^\infty$ converges to $\nu|_T$ in $\Dr(M|_T)$. 
\end{proposition}
\begin{proof}
 Fix a choice of $I\subset [n]\sm T$ of complementary rank such that $T\cup I$ spans $M$. Consider that $\d(\nu_i,\nu)\geq \d(\pi_{T\cup I}(\nu_i),\imath(\nu|_T))$ implies that the sequence $\set{\pi_{T\cup I}(\nu_i)}_{i=1}^\infty$ converges to $\imath(\nu|_T)$. Since $\imath_{I}$ is an isometry, convergence of $\set{\pi_{T\cup I}(\nu_i)}_{i=1}^\infty$ implies convergence of $\set{\imath\inv(\pi_{T\cup I}(\nu_i|_T))}_{i=1}^\infty$ to $\nu_T$. In particular we have that 
 $$\imath\inv(\pi_{T\cup I}(\nu_i|_T))=\nu_i|_T$$
 and the result follows.
\end{proof}
}